% =====================================================================
%  SkillSeek AI — Guide de démarrage
%  Compilation : pdflatex guide.tex
% =====================================================================
\documentclass[11pt,a4paper]{article}

\usepackage[utf8]{inputenc}
\usepackage[T1]{fontenc}
% Typographie française si babel-french est disponible (MiKTeX l'installe
% à la volée) ; sinon on se contente de l'espacement, et le document compile
% quand même. Un guide de démarrage qui ne compile pas serait une ironie.
\IfFileExists{french.ldf}{\usepackage[french]{babel}}{\frenchspacing}
\usepackage[margin=2.4cm]{geometry}
\usepackage{xcolor}
\usepackage{listings}
\usepackage{enumitem}
\usepackage{array}
\usepackage{booktabs}
\usepackage{titlesec}
\usepackage[colorlinks=true,linkcolor=bleu,urlcolor=bleu]{hyperref}

\definecolor{bleu}{HTML}{1F4E79}
\definecolor{gris}{HTML}{595959}
\definecolor{fondcode}{HTML}{F2F4F7}

% Titres en bleu, sobres
\titleformat{\section}{\normalfont\Large\bfseries\color{bleu}}{\thesection.}{0.6em}{}
\titleformat{\subsection}{\normalfont\large\bfseries\color{bleu}}{\thesubsection}{0.6em}{}
\titlespacing*{\section}{0pt}{1.4em}{0.6em}

% Blocs de commandes : fond gris, chasse fixe, pas de coloration inutile
\lstset{
  basicstyle=\ttfamily\small,
  backgroundcolor=\color{fondcode},
  frame=none,
  framexleftmargin=8pt,
  framexrightmargin=8pt,
  xleftmargin=8pt,
  aboveskip=10pt,
  belowskip=10pt,
  breaklines=true,
  columns=fullflexible,
  keepspaces=true,   % sans quoi l'alignement des commentaires est écrasé
  literate={é}{{\'e}}1 {è}{{\`e}}1 {ê}{{\^e}}1 {à}{{\`a}}1 {ç}{{\c c}}1 {É}{{\'E}}1,
}

\setlist[itemize]{leftmargin=1.2em,itemsep=2pt,topsep=3pt}
\renewcommand{\arraystretch}{1.3}

\begin{document}

% ---------------------------------------------------------------- Titre
\begin{center}
  {\Huge\bfseries\color{bleu} SkillSeek AI}\\[0.5em]
  {\Large\color{gris} Guide de démarrage}\\[1.2em]
  {\small\itshape\color{gris}
   Plateforme d'aide au recrutement par analyse automatique des CV}\\[2em]
  {\small Aymen Benrbib \quad$\cdot$\quad Stage d'études — BC Skills
   \quad$\cdot$\quad ESI Rabat}
\end{center}
\vspace{1.5em}

\noindent
SkillSeek AI analyse les CV déposés sur une offre, en tire un profil
structuré et attribue une note de correspondance dont elle donne le détail.
La plateforme classe et explique ; elle ne décide pas. Tout tourne en
conteneurs Docker : ni Python ni Node à installer.

% ---------------------------------------------------------------- 1
\section{Prérequis}

\begin{itemize}
  \item \textbf{Docker Desktop} --- \url{https://docker.com/products/docker-desktop}
  \item \textbf{Git} --- \url{https://git-scm.com}
\end{itemize}

\noindent
Rien d'autre. \textbf{Docker Desktop doit être démarré} avant les commandes
qui suivent : c'est la cause d'échec la plus fréquente.

% ---------------------------------------------------------------- 2
\section{Démarrage}

Dans PowerShell (Windows) ou un terminal (macOS, Linux). Une commande à la
fois : chacune doit se terminer avant la suivante.

\subsection*{Récupérer le code}
\begin{lstlisting}
git clone -b livraison-bc-skills https://github.com/AyMB1303/skillseek-ai.git
cd skillseek-ai
\end{lstlisting}
\noindent
\textcolor{gris}{\small\itshape
\texttt{livraison-bc-skills} est la branche de livraison : son contenu ne
bougera pas. Les évolutions ultérieures arriveront sur \texttt{main}.}

\subsection*{Créer le fichier de configuration}
Le dépôt ne contient aucun mot de passe. Le modèle fourni suffit en local.
\begin{lstlisting}
copy .env.example .env     # Windows
cp .env.example .env       # macOS / Linux
\end{lstlisting}

\subsection*{Lever la plateforme}
\begin{lstlisting}
docker compose up -d --build
\end{lstlisting}
La première exécution construit les images : de cinq à quinze minutes. Les
fois suivantes, quelques secondes.

\subsection*{Préparer la base de données}
Dans cet ordre : le schéma, les rôles, puis un jeu de démonstration.
\begin{lstlisting}
docker compose exec backend flask db upgrade
docker compose exec backend flask seed
docker compose exec backend flask demo --reset
\end{lstlisting}

% ---------------------------------------------------------------- 3
\section{Vérifier}

\begin{center}
\begin{tabular}{@{}>{\ttfamily}p{6.2cm}p{8.4cm}@{}}
\toprule
\textrm{\textbf{Adresse}} & \textbf{Ce que vous devez voir}\\
\midrule
http://localhost:3000 & L'interface, page de connexion\\
http://localhost:5000/api/health & Une réponse JSON de l'API\\
\bottomrule
\end{tabular}
\end{center}

\noindent Si les deux répondent, l'installation est terminée.

% ---------------------------------------------------------------- 4
\section{Comptes de démonstration}

\begin{center}
\begin{tabular}{@{}p{3.2cm}>{\ttfamily}p{6.2cm}>{\ttfamily}p{3.4cm}@{}}
\toprule
\textbf{Rôle} & \textrm{\textbf{Identifiant}} & \textrm{\textbf{Mot de passe}}\\
\midrule
Candidat        & y.tazi@example.ma      & Demo@1234\\
Recruteur       & s.lamrani@bcskills.ma  & Demo@1234\\
Administrateur  & admin@skillseek.local  & Admin@1234\\
\bottomrule
\end{tabular}
\end{center}

% ---------------------------------------------------------------- 5
\section{Le parcours à regarder}

\textbf{Candidat} --- ouvrez une offre (ses critères sont annoncés avant
tout dépôt), déposez un CV au format PDF. L'analyse est immédiate.

\medskip
\noindent
\textbf{Recruteur} --- dans « Candidatures », ouvrez « Détails » sur un
dossier. C'est là que se joue l'essentiel :

\begin{itemize}
  \item le détail du calcul : d'où viennent les points, composante par composante ;
  \item les réserves et les critères éliminatoires, avec leur motif exact ;
  \item le profil reconstitué, qui indique aussi ce qui ne figure pas dans le CV ;
  \item le CV lui-même, affiché à côté du profil, pour vérification directe ;
  \item l'avis du modèle appris, borné à $\pm 8$ points : il ajuste la note,
        il ne la décide pas.
\end{itemize}

% ---------------------------------------------------------------- 6
\section{Commandes courantes}

\begin{center}
\begin{tabular}{@{}p{5.2cm}>{\ttfamily\small}p{9.4cm}@{}}
\toprule
\textbf{Besoin} & \textrm{\textbf{Commande}}\\
\midrule
Arrêter          & docker compose down\\
Relancer         & docker compose up -d\\
Voir les journaux& docker compose logs -f backend frontend\\
Données neuves   & docker compose exec backend flask demo --reset\\
Lancer les tests & docker compose exec backend pytest -q\\
Tout effacer     & docker compose down -v\\
\bottomrule
\end{tabular}
\end{center}

\noindent
\textcolor{gris}{\small\itshape
« docker compose down -v » supprime aussi la base : il faut ensuite reprendre
la préparation des données.}

% ---------------------------------------------------------------- 7
\section{En cas de problème}

\begin{center}
\begin{tabular}{@{}p{5.6cm}p{9.0cm}@{}}
\toprule
\textbf{Symptôme} & \textbf{Cause et remède}\\
\midrule
\texttt{cannot connect to the Docker daemon}
 & Docker Desktop n'est pas démarré.\\
\texttt{port is already allocated}
 & Un autre programme occupe le port 3000, 5000 ou 5432.\\
La page 3000 tarde
 & Le serveur compile à la demande : la première ouverture peut prendre une minute.\\
\texttt{relation does not exist}
 & La préparation de la base n'a pas été faite, ou pas dans l'ordre.\\
Assistant RH sommaire
 & Le modèle de langue local est facultatif et absent. La plateforme bascule sur une rédaction déterministe : c'est prévu, ce n'est pas une panne.\\
\bottomrule
\end{tabular}
\end{center}

% ---------------------------------------------------------------- 8
\section{Un mot sur la méthode}

\textbf{Le système propose, l'humain tranche.} Aucune décision n'est
automatique.

\medskip
\noindent
\textbf{Ce qui n'a pas pu être mesuré ne s'écrit jamais comme une mesure.}
Une rubrique absente du CV est annoncée comme telle plutôt que masquée ; une
expérience illisible place la candidature en réserve au lieu de l'écarter ;
du texte posé hors du cadre de la page --- invisible à la lecture --- est
écarté de l'analyse et signalé.

\end{document}
